% from ega4-17-fr.tex
\oldpage[IV-4]{60}
\begin{enumerate}
  \item[(i)] On sait (16.3.9) que si $f$ est localement de type fini, $\Omega_f^1$ est un $\sh{O}_X$-Module de type fini. Pour prouver que, dans tous les cas, il est localement projectif, on peut se borner, en vertu de (17.1.6), au cas où $Y=\Spec(A)$ et $X=\Spec(B)$ sont affines, et cela résulte de l'hypothèse sur $f$ et de (0, 20.4.9 et 0, 19.2.1).
  \item[(ii)] Ici encore, on peut se borner au cas où $X$, $Y$ et $Z$ sont affines (17.1.6) et la conclusion résulte dans ce cas de l'interprétation des Modules figurant dans la suite (17.2.3.1) et de (0, 20.5.7).
\end{enumerate}

\begin{corollary}[17.2.4]
\label{IV.17.2.4}
Si $f:X\to Y$ est un morphisme formellement étale, alors, pour tout morphisme $g:Y\to Z$, l'homomorphisme canonique de $\sh{O}_X$-Modules
\[
  f^*(\Omega_{Y/Z}^1)\longrightarrow\Omega_{X/Z}^1
\]
est bijectif.
\end{corollary}

Cela résulte de l'exactitude de la suite (17.2.3.1) et du fait que l'on a alors $\Omega_{X/Y}^1=0$ (17.2.1).

\begin{proposition}[17.2.5]
\label{IV.17.2.5}
Soient $f:X\to Y$ un morphisme, $X'$ un sous-préschéma de $X$ tel que le morphisme composé $X'\xrightarrow{j}X\to Y$ (où $j$ est l'injection canonique) soit formellement lisse. Alors la suite de $\sh{O}_{X'}$-Modules (16.4.21)
\[
\label{IV.17.2.5.1}
  0\longrightarrow\mathcal{N}_{X'/X}\longrightarrow\Omega_{X/Y}^1\otimes_{\sh{O}_X}\sh{O}_{X'}\longrightarrow\Omega_{X'/Y}^1\longrightarrow0
  \tag{17.2.5.1}
\]
est exacte; en outre, pour tout $x\in X$, il existe un voisinage ouvert $U$ de $x$ tel que les restrictions à $U$ des homomorphismes de (17.2.5.1) forment une suite exacte et scindée.
\end{proposition}

Toujours en vertu de (17.1.6), on peut se borner au cas où $Y=\Spec(A)$ et $X=\Spec(B)$ sont affines, et $X'=\Spec(B/\mathfrak{J})$, où $\mathfrak{J}$ est un idéal de $B$. Alors le faisceau conormal $\mathcal{N}_{X'/X}$ correspond au $B$-module $\mathfrak{J}/\mathfrak{J}^2$ (16.1.3), et la conclusion découle de (0, 20.5.14).

\begin{proposition}[17.2.6]
\label{IV.17.2.6}
Soient $X$, $Y$ deux préschémas, $f:X\to Y$ un morphisme localement de type fini. Les conditions suivantes sont équivalentes:
\begin{enumerate}
  \item[a)] $f$ est un monomorphisme.
  \item[b)] $f$ est radiciel et formellement non ramifié.
  \item[c)] Pour tout $y\in Y$, la fibre $f^{-1}(y)$ est vide ou $k(y)$-isomorphe à $\Spec(k(y))$ (autrement dit, est réduite à un seul point $z$ tel que $k(y)\to\sh{O}_z/\mathfrak{m}_y\sh{O}_z$ soit un isomorphisme).
\end{enumerate}
\end{proposition}

Le fait que a) entraîne c) résulte de (8.11.5.1). Il est clair que c) entraîne que $f$ est radiciel; montrons qu'il résulte aussi de c) que $\Omega_{X/Y}^1=0$, ce qui prouvera que c) entraîne b) (17.2.1). Notons que le $\sh{O}_X$-Module $\Omega_{X/Y}^1$ est quasi-cohérent de type fini (16.3.9). Il résulte donc de (I, 9.1.13.1) que, pour que $(\Omega_{X/Y}^1)_x=0$, il faut et il suffit que si l'on pose $Y_1=\Spec(k(y))$, $X_1=f^{-1}(y)=X\times_Y Y_1$, on ait $(\Omega_{X_1/Y_1}^1)_x=0$; mais comme le morphisme $f_1:X_1\to Y_1$ déduit de $f$ est formellement non ramifié en vertu de l'hypothèse c) (17.1.3), la conclusion résulte de (17.2.1). Prouvons enfin que b) entraîne a); pour cela, considérons le morphisme diagonal $g=\Delta_f:X\to X\times_Y X$; puisque $f$ est radiciel, $g$ est surjectif (I, 8.7.1); d'autre part, $\Omega_{X/Y}^1$ est par définition le faisceau conormal $\mathop{Gr}\nolimits_1(g)$ de l'immersion $g$ (16.3.1), et dire que $f$ est formellement non ramifié signifie donc que

